%% =====================================================================
%%  RESEARCH STATEMENT
%% =====================================================================
%%  \cite{key} prints the label your CV gives that publication, e.g.
%%  [J5], and links to it. Several keys: \cite{key1,key2} -> [J5, C3].
%%  Figures work as usual: \begin{figure} ... \caption{...} \end{figure}.
%%  Replace the TikZ diagram below with \includegraphics if you prefer.
%% =====================================================================
\begin{applicationdocument}[
  page limit = 4,
  subtitle   = {Prepared for the position of \PositionSummary},
]{Research Statement}

\tip[What committees look for]{The research statement is where a
committee judges your \emph{future} program: is it important, original,
feasible and fundable, and does it fit here? Check the posting for length.
Two to four pages is common; research-intensive posts sometimes allow five.
Open with a vision a non-specialist can follow. Present past work as the
foundation for what comes next, not as a history. Give most of the space to
specific aims, first student projects and a funding plan.}

Lakes and reservoirs supply drinking water to millions of Canadians, yet
the harmful algal blooms that threaten them are becoming more frequent,
more intense and harder to predict. Water managers need to know where and
when a bloom will appear early enough to act: to adjust treatment, issue
advisories or schedule sampling. My research asks a question at the heart
of environmental data science: \textbf{how can we make environmental
forecasts reliable when the future no longer resembles the past?} I answer
it by combining process-based models, machine learning and
community-collected data, and by building forecasts with the people who use
them. At \Institution, my lab will develop forecasting methods that stay
trustworthy under climate change and draw on community monitoring under
community control. It will also test whether better forecasts lead to
better decisions.

\section{Research to Date}

\textbf{Hybrid models that know some physics.} Purely data-driven models
fail when conditions move outside their training data. Process-based
models, for their part, are expensive to calibrate and miss biological feedbacks. In my
doctoral work I designed process-guided neural networks that build the
physics of lake mixing and heat transfer into their architecture
\cite{ellis2025hybrid,ellis2022physics}. Across 18 lakes, these hybrid models
forecast bloom onset ten days ahead with 35\% lower error than either model
family alone. They kept that advantage in years warmer than any in their
training data. I also developed training strategies for the sparse,
irregular observations that are the norm in lake monitoring
\cite{ellis2023sparse}.

\textbf{Uncertainty that people can use.} A forecast is only useful if its
uncertainty is honest. I showed that standard ensemble methods are
overconfident for lake water quality. I then developed a calibration
method that brings 90\% prediction intervals to their stated coverage in
every season \cite{ellis2024uncertainty}. With my trainee Haruto Nakamura, I
turned uncertainty into action. Sampling schedules that send field crews
where forecasts are least certain detect the same number of blooms with
30\% fewer site visits \cite{nakamura2024active}. Most recently, we developed
tests that flag when a model has drifted outside the conditions it knows,
for example during heatwaves and storms \cite{ellis2026shift}.

\textbf{From research code to operational forecasts.} As a postdoctoral
fellow I led the development of BloomCast, an open-source system that issues
daily probabilistic forecasts for 42 lakes \cite{ellis2025bloomcast}. Two
provincial agencies use it to plan sampling, which cut the median time to
detect a bloom from nine days to three. With lake associations and a First
Nation's water monitoring program, we showed that volunteer measurements
improve forecasts at sparsely monitored lakes when their quality is modelled
explicitly \cite{martin2025community}. These partnerships taught me that the
hardest open problems are not only technical. They also concern who
controls the data and how forecasts change decisions. Those questions shape
the program I propose below.

\section{Research Program at \Institution}

\tip[Make the future concrete]{Committees weigh feasibility as heavily as
ambition. For each aim, give the question, the approach, the first project
a student could start in Year~1, and the expected outcome. A one-glance
figure helps reviewers who skim. Connect each aim to the work above so
your track record vouches for it.}

My program for the next five to seven years has three connected aims
(Figure~\ref{fig:program}). Each builds directly on my published work, and
each begins with a project a student can start in my first year.

\begin{figure}[htb]
  \centering
  \begin{tikzpicture}[
      font=\sffamily\footnotesize,
      aim/.style={draw=accent, line width=0.7pt, fill=accent!5,
        rounded corners=3pt, text width=0.27\linewidth, inner sep=6pt,
        align=left, anchor=north, minimum height=3.15cm,
        execute at begin node={\hyphenpenalty=10000\exhyphenpenalty=10000}},
      goal/.style={fill=accent, text=white, rounded corners=3pt,
        inner sep=7pt, align=center, font=\sffamily\small\bfseries},
      arrow/.style={-{Stealth[length=5pt]}, accent, line width=0.8pt},
    ]
    \node[aim] (a1) at (0,0) {%
      {\color{accent}\bfseries AIM 1}\\
      \textbf{Forecasting under change}\\[3pt]
      Models that detect and adapt to conditions they have never seen.\\[3pt]
      {\color{black!60}First project: shift-aware hybrid models (Ph.D.)}};
    \node[aim] (a2) at (0.333\linewidth,0) {%
      {\color{accent}\bfseries AIM 2}\\
      \textbf{Community data, community control}\\[3pt]
      Forecasts that learn from volunteer and Indigenous-led monitoring.\\[3pt]
      {\color{black!60}First project: federated training with partners (M.Sc.)}};
    \node[aim] (a3) at (0.666\linewidth,0) {%
      {\color{accent}\bfseries AIM 3}\\
      \textbf{Decision-ready forecasts}\\[3pt]
      Field experiments on how forecasts change real decisions.\\[3pt]
      {\color{black!60}First project: trials with conservation authorities (M.Sc.)}};
    \node[goal, minimum width=0.946\linewidth] (goal)
      at ([yshift=-0.85cm]a2.south)
      {Reliable, equitable freshwater forecasts that people act on};
    \foreach \n in {a1,a2,a3}
      \draw[arrow] (\n.south) -- (\n.south |- goal.north);
  \end{tikzpicture}
  \caption{My research program at a glance. Each aim builds on published
    work and starts with a student project in Year~1.}
  \label{fig:program}
\end{figure}

\subsection{Aim 1: Forecasting under change}
\emph{How can a forecasting model recognize conditions it has never seen,
and adapt to them?} I will extend my shift-detection tests
\cite{ellis2026shift} into models that adapt as new observations arrive,
using process knowledge to keep their extrapolations physically plausible.
Two resources will drive this work: the LakeBench benchmark of 120 lakes
\cite{patel2026benchmark}, and the Westbrook Centre for Freshwater Futures'
30-year record for the Westbrook River watershed. A Ph.D.\ student will
begin with shift-aware hybrid models. The outcome will be forecasts that
know when they do not know, released through BloomCast.

\subsection{Aim 2: Community data under community control}
\emph{How can forecasts learn from volunteer and Indigenous-led monitoring
while the data stay with the people who collect them?} I will develop
federated training methods, which let models learn from data that never
leave the community's own systems. I will co-design them with partner
communities under data agreements based on the First Nations principles of
OCAP\textsuperscript{\textregistered}. A master's student co-supervised with
a community partner will lead the first project. The outcome will be
forecasts for lakes that no agency monitors, and a data-governance model
other projects can adopt.

\subsection{Aim 3: Decision-ready forecasts}
\emph{Do better forecasts lead to better decisions, and which forecast
products help most?} Building on our pilot experiment with water managers
\cite{ellis2026decisions}, I will run randomized field trials with
conservation authorities. I will also measure the value of forecasts in
decision-theoretic terms: how much they improve the choices managers make.
A master's student and two undergraduate researchers will start with
sampling and advisory decisions. The outcome will be direct evidence of
what makes a forecast useful, and products co-designed with the people who
rely on them.

In the longer term, these aims position my lab to coordinate a regional
freshwater forecasting network. The same approach will extend from algal
blooms to drought, road-salt contamination and other emerging threats to
water quality.

\section{Funding Plan}

\begin{datatable}{@{} l X X @{}}
  \colhead{When} & \colhead{Source} & \colhead{Supports}\\
  \midrule\endhead
  Year 1   & NSERC Discovery Grant, with Discovery Launch Supplement
           & Aims 1--2: two graduate students, computing\\
  Years 1--2 & NSERC Alliance, with provincial agencies and conservation
               authorities (two partners have committed)
           & Aim 3: field trials, partner-site sampling\\
  Year 2   & CFI John R.\ Evans Leaders Fund, with provincial matching
           & Real-time sensor network on Westbrook River lakes\\
  Ongoing  & Mitacs Accelerate; NSERC Undergraduate Student Research Awards
           & Internships with partners; paid undergraduate research\\
  \bottomrule
\end{datatable}

My record supports this plan. I contributed the forecasting work package to
a successful \$480{,}000 NSERC Alliance grant, I hold a \$25{,}000 seed grant
as principal investigator, and I was ranked third of 41 applicants for my
NSERC Postdoctoral Fellowship.

\section{Training, Collaboration and Fit}

By Year~3 I expect a lab of two Ph.D.\ students, two master's students and
three or four undergraduate researchers. They will learn hybrid modelling,
open and reproducible software, and partnered research. I will fund every
position, including undergraduate ones, and set shared expectations through
a written lab charter. \Institution is an ideal home for this program. The
Westbrook Centre for Freshwater Futures offers field sites and long-term
data. Dr.~Priya Raman's work on watershed hydrology would let us link river
inflows to lake forecasts in Aim~1. Dr.~Tom\'as Ortega's work in aquatic
ecology complements Aim~3. The \Faculty's data science computing cluster
meets my computing needs. I will keep my existing collaborations with the
University of Eastvale, the Nordic Lake Observatory and the agencies that
use BloomCast. Together, these give my students a path from methods to
real-world use, and give \Institution a forecasting program built to serve
its region.

\end{applicationdocument}
